\chapter{Arbitrary-Order Production and Fixed-Prime Sampling}
\label{ch:universality}

\section{Continuous density-of-states interpolation}

Let
\begin{equation}
\alpha:=\frac1\sigma,
\qquad
m_h\simeq\Lambda p^{\sigma h},
\qquad
N_h=p^h-1.
\label{eq:univ-alpha}
\end{equation}
Interpolating the endpoint counting function gives
\begin{equation}
N_{\mathrm c}(m)=\left(\frac m\Lambda\right)^\alpha-1,
\qquad
\frac{\dd N_{\mathrm c}}{\dd m}
=\frac1{\sigma\Lambda^\alpha}m^{\alpha-1}.
\label{eq:univ-cont-count}
\end{equation}
This defines a continuous model.  It is exact as an interpolation of the endpoint law.  \Cref{sec:univ-fixed-prime} shows why it is not an ordinary pointwise limit of a fixed-prime tower.

At a prescribed linear crossing, take
\begin{equation}
\nu(m)=Am,
\qquad
f(k,m)=\eta\exp\left[-\pi\left(\frac mA+\frac{k^2}{Am}\right)\right].
\label{eq:univ-distribution}
\end{equation}
The Gaussian momentum integral gives the number spectrum
\begin{equation}
\boxed{
\frac{\dd n}{\dd m}
=\frac{\eta A^{3/2}}
{8\pi^3\sigma\Lambda^\alpha}
m^{\alpha+1/2}\ee^{-\pi m/A}.}
\label{eq:univ-number-spectrum}
\end{equation}
Multiplication by $m$ gives the captured rest-energy spectrum
\begin{equation}
\boxed{
\frac{\dd\rho_R}{\dd m}
=\frac{\eta A^{3/2}}
{8\pi^3\sigma\Lambda^\alpha}
m^{\alpha+3/2}\ee^{-\pi m/A}.}
\label{eq:univ-rest-spectrum}
\end{equation}

\section{Exact arbitrary-order constants}

\begin{theorem}[Continuous arbitrary-order production]
For $\sigma>0$, the continuous model \eqref{eq:univ-cont-count}--\eqref{eq:univ-distribution} has
\begin{align}
n_\star
&=C_N(\sigma)\eta\frac{A^{\alpha+3}}{\Lambda^\alpha},
&C_N(\sigma)
&=\frac{\Gamma(\alpha+\tfrac32)}
{8\sigma\pi^{\alpha+9/2}},
\label{eq:univ-CN}\\
\rho_{R,\star}
&=C_R(\sigma)\eta\frac{A^{\alpha+4}}{\Lambda^\alpha},
&C_R(\sigma)
&=\frac{\Gamma(\alpha+\tfrac52)}
{8\sigma\pi^{\alpha+11/2}},
\label{eq:univ-CR}\\
\rho_\star
&=C_E(\sigma)\eta\frac{A^{\alpha+4}}{\Lambda^\alpha},
&C_E(\sigma)
&=\frac{\Gamma(\alpha+4)B(\tfrac32,\alpha+2)}
{4\sigma\pi^{\alpha+6}},
\label{eq:univ-CE}\\
p_\star
&=C_P(\sigma)\eta\frac{A^{\alpha+4}}{\Lambda^\alpha},
&C_P(\sigma)
&=\frac{C_E(\sigma)}{2\alpha+7}.
\label{eq:univ-CP}
\end{align}
Consequently,
\begin{equation}
\boxed{
w_\star(\sigma)=\frac1{7+2/\sigma},}
\label{eq:univ-w}
\end{equation}
\begin{equation}
\boxed{
f_R(\sigma)=\frac{\rho_{R,\star}}{\rho_\star}
=\frac{C_R(\sigma)}{C_E(\sigma)},}
\label{eq:univ-rest-frac}
\end{equation}
and
\begin{equation}
\boxed{
\frac{\rho_{R,\star}}{n_\star}
=m_{\mathrm{peak}}^{(R)}
=\frac A\pi\left(\frac1\sigma+\frac32\right).}
\label{eq:univ-peak}
\end{equation}
\end{theorem}

\begin{proof}
Equations \eqref{eq:univ-CN} and \eqref{eq:univ-CR} follow by integrating \eqref{eq:univ-number-spectrum} and \eqref{eq:univ-rest-spectrum} with the gamma integral.

For total energy and pressure, set $k=m\sinh\theta$.  The production exponent becomes
\begin{equation}
\pi\frac mA\cosh^2\theta.
\end{equation}
After the mass integral, the energy angular factor is
\begin{equation}
I_E=\int_0^\infty
\sinh^2\theta\cosh^{-2\alpha-6}\theta\,\dd\theta
=\frac12B\left(\frac32,\alpha+2\right).
\end{equation}
The pressure angular factor is
\begin{equation}
I_P=\frac13\int_0^\infty
\sinh^4\theta\cosh^{-2\alpha-8}\theta\,\dd\theta
=\frac16B\left(\frac52,\alpha+2\right).
\end{equation}
Their ratio is
\begin{equation}
\frac{I_P}{I_E}
=\frac13\frac{B(\tfrac52,\alpha+2)}{B(\tfrac32,\alpha+2)}
=\frac1{2\alpha+7},
\end{equation}
which proves \eqref{eq:univ-w}.  The gamma recurrence gives the mean rest mass, and differentiating \eqref{eq:univ-rest-spectrum} gives the same value for its peak.
\end{proof}

At $\sigma=1$, these formulas reduce to
\begin{equation}
w_\star=\frac19,
\qquad
f_R=\frac{525\pi}{2048},
\qquad
m_{\mathrm{peak}}^{(R)}=\frac{5A}{2\pi}.
\end{equation}

\begin{figure}[htbp]
\centering
\includegraphics[width=0.98\linewidth]{general_sigma_universality.pdf}
\caption{The continuous arbitrary-order family.  Increasing $\sigma$ raises the production equation of state while lowering the captured rest fraction and the selected mass in units of $A$.}
\label{fig:general-sigma}
\end{figure}

\section{Prime independence of the interpolation}

The prime $p$ cancels from \eqref{eq:univ-cont-count} because
\begin{equation}
p^h=\left(\frac{m_h}{\Lambda}\right)^{1/\sigma}.
\end{equation}
This gives a continuous universality class labeled by $\sigma$.  Prime independence at this stage does not mean that the discrete tower forgets $p$.  The spacing in $\log m$ is $\sigma\log p$, which remains finite when $\Lambda/A$ is reduced.

\section{Fixed-prime Mellin correction}
\label{sec:univ-fixed-prime}

Let $F(y)$ be a single-state dimensionless observable with $y=m/A$, and define
\begin{equation}
x=\frac\Lambda A,
\qquad
c_p=\frac{p-1}{p},
\qquad
a=\sigma\log p.
\end{equation}
Ignoring a lower-end correction that vanishes as $x\to0$ for the production kernels, extend the level sum bilaterally:
\begin{equation}
S_F(x)=c_p\sum_{h\in\mathbb Z}p^hF(xp^{\sigma h}).
\label{eq:univ-discrete-sum}
\end{equation}
Let
\begin{equation}
\Mellin F(s)=\int_0^\infty y^{s-1}F(y)\,\dd y.
\end{equation}

\begin{theorem}[Log-periodic fixed-prime decomposition]
Suppose $F(y)=\order(y^\beta)$ as $y\to0$ for some $\beta>0$, decays sufficiently rapidly at infinity, and
\begin{equation}
g_x(t)=c_p\ee^{t\log p}F(x\ee^{at})
\end{equation}
together with its Fourier transform satisfies the integrability and summability hypotheses of Poisson summation.  A sufficient stronger condition is that $g_x$ is Schwartz.  Then
\begin{equation}
\boxed{
S_F(x)
=\frac{c_p}{a}x^{-\alpha}
\sum_{k\in\mathbb Z}
x^{\ii\omega_k}
\Mellin F(\alpha-\ii\omega_k),
\qquad
\omega_k=\frac{2\pi k}{a}.}
\label{eq:univ-mellin-poisson}
\end{equation}
\end{theorem}

\begin{proof}
Apply Poisson summation to
\begin{equation}
g(t)=c_p\ee^{t\log p}F(x\ee^{at}).
\end{equation}
In the Fourier transform, substitute $y=x\ee^{at}$.  The factor $\ee^{-2\pi\ii kt}$ becomes $(y/x)^{-\ii\omega_k}$, and the remaining integral is the Mellin transform in \eqref{eq:univ-mellin-poisson}.  The small-$y$ hypothesis also bounds the omitted physical levels $h\le0$.  Their absolute contribution is $\order(x^\beta)$ and is negligible relative to the $x^{-\alpha}$ leading abundance.
\end{proof}

The $k=0$ term has coefficient
\begin{equation}
\frac{c_p}{\sigma\log p}x^{-\alpha}\Mellin F(\alpha).
\label{eq:univ-zero-harmonic}
\end{equation}
The derivative interpolation \eqref{eq:univ-cont-count} instead has coefficient
\begin{equation}
\frac1\sigma x^{-\alpha}\Mellin F(\alpha).
\end{equation}
Their mean amplitude ratio is
\begin{equation}
\boxed{
\frac{c_p}{\log p}
=\frac{p-1}{p\log p}.}
\label{eq:univ-amplitude-ratio}
\end{equation}
It approaches one only in a formal dense-branching limit $p\to1$, not along the physical prime sequence.  The $k\ne0$ terms add a periodic dependence on $\log x$.

\begin{corollary}[No fixed-prime convergence to the derivative interpolation]
After the natural $x^\alpha$ rescaling, reducing $\Lambda/A$ at fixed $p$ translates the sampling grid in $\log m$ but does not shrink its spacing.  The rescaled discrete normalization therefore retains a prime-dependent mean offset and log-periodic corrections rather than converging to the derivative interpolation.  Ratios of observables can nevertheless be much more accurate because common amplitude factors cancel.
\end{corollary}

\clearpage
\begin{figure}[htbp]
\centering
\includegraphics[width=0.82\linewidth]{fixed_prime_log_sampling.pdf}
\caption{Relative discrepancy between the discrete total production energy and the derivative interpolation at $\sigma=1$.  Decreasing $\Lambda/A$ does not remove the fixed logarithmic spacing.  The nearly constant dyadic offset and the larger prime-dependent oscillations agree with the Mellin decomposition.}
\label{fig:fixed-prime-sampling}
\end{figure}

\section{Correction to the interpretation of the continuum formulas}

The continuous constants remain exact for the density-of-states model defined by \eqref{eq:univ-cont-count}.  They also organize robust ratios in the dense mass band.  They should not be described as a convergent fixed-prime limit without an additional averaging operation, a refinement in logarithmic spacing, or a theorem that suppresses the nonzero Mellin harmonics.

This distinction matters most for absolute abundance.  The cosmological benchmark uses the direct discrete tower calculation when imposing its energy projection, so its discrete normalization is not replaced by the continuous coefficient.  The estimate explains why the reported dyadic ratios can be close to the continuous values even though the absolute energy coefficient differs substantially.
